\subsection{Positive-dimensional henselizations over effective $C_1$ fields}
\label{C:C1-henselizations-subsection}

In this subsection $k$ is an infinite countable $C_1$ field
of characteristic $p>2$, with a computable surjective
enumeration, exact decidable arithmetic and equality,
and univariate factorization over every explicitly
presented finite tower of field extensions of $k$.
These are the finite-input hypotheses of
Theorem~\ref{C:artinian_effective_imperfect_C1};
perfectness is not assumed.

\begin{lemma}[Proved: exact arithmetic in affine local henselizations over $k$]
\label{C:C1-henselization-access}
Let $S=k[z_1,\ldots,z_s]/I$ be a supplied finite
presentation, let $\kappa/k$ be a supplied finite
tower of field extensions, and let
$\rho_S:S\twoheadrightarrow\kappa$ be a supplied
surjective $k$-algebra map.  Put
$\mathfrak m=\ker\rho_S$ and $R=(S_{\mathfrak m})^h$.
The ring $R$ has exact finite names with decidable
validity and equality, computable arithmetic,
residue and unit inversion, and computable lifts
of every element of $\kappa$.
An ordinary terminating algorithm computes generators
of the full actual kernel of every finite matrix
over $R$.  No regularity or reducedness of $S$
is required.
\end{lemma}
\begin{proof}
Finite polynomial computations over $k$ use its
exact field arithmetic.  The Gr\"obner construction
and its module and syzygy variants require only
field operations and decidable equality, and
terminate by the ascending condition on monomial
submodules.  The explicit constructions with
tracked coefficients are given in
\cite[Sections 4.3--4.4 and 5.1--5.3,
pp.\ 240--249]{C-BuchbergerSystems2001}.
Thus ideal membership and polynomial matrix
kernels are computable over $k$.
For a presented algebra
$T=k[w_1,\ldots,w_a]/(F_1,\ldots,F_b)$
and a matrix $A:T^u\to T^v$, choose a polynomial
lift $\widetilde A$ and append the columns
$-F_j e_l$ for $1\leq j\leq b$, $1\leq l\leq v$.
Project generators of the polynomial syzygy
module to their first $u$ coordinates and reduce
them in $T$.  They generate $\ker A$, since
$\widetilde A x$ is zero in $T^v$ exactly when
it is a polynomial combination of the $F_j e_l$.
Taking multiplication by one element computes
its annihilator.  Inverting finitely many elements
is represented by adjoining variables and equations
$u_jg_j-1$, so the same algorithms apply to all
finite localizations used below.

We name elements using pointed standard \'etale
neighborhoods.  Invert finitely many elements of
$S$ with nonzero $\rho_S$-image, and let $S'$ be
that presented localization.  Choose a monic
$f\in S'[U]$, $h\in S'[U]$, and $a\in\kappa$
such that
\[
 \overline f(a)=0,\qquad
 \overline{f'}(a)\,\overline h(a)\ne0.
                                                               \tag{C.3.CH1}
\]
Take
$T=(S'[U]/(f))_{f'h}$ with residue map
$\rho:T\to\kappa$ sending $U$ to $a$,
and take an element $z\in T$.
These are finite data with decidable validity.
The map is standard \'etale, and $\rho$ is
surjective because its restriction to $S$ is.
Write $\mathfrak q=\ker\rho$.
The chosen simple root gives a canonical map
$T\to R$ through this point.

All elements of $R$ have such names.  The
henselization is the filtered colimit of pointed
\'etale neighborhoods with unchanged residue,
and near the point each \'etale neighborhood is
standard \'etale
\cite[Lemma 10.155.1 and Proposition 10.144.4,
Tags 0BSK and 00UE]{C-StacksHenselianEffective}.
Finite presentation descends its finitely many
coefficients and denominators from $S_{\mathfrak m}$
to one such $S'$.
The pointed local map identifies
$(T_{\mathfrak q})^h$ with $R$.
Consequently $T\to R$ is flat and
$T_{\mathfrak q}\to R$ is faithfully flat
\cite[Section 10.155 and Lemma 15.46.1,
Tags 0BSK and 07QL]{C-StacksHenselianEffective}.
For every named element,
\[
 z=0\text{ in }R
 \quad\Longleftrightarrow\quad
 \operatorname{Ann}_T(z)\not\subseteq\mathfrak q.
                                                               \tag{C.3.CH2}
\]
Indeed faithful flatness reflects zero in
$T_{\mathfrak q}$, and localization kills $z$
exactly when an element outside $\mathfrak q$
annihilates it.  Compute finitely many annihilator
generators by the first paragraph and test their
residues.  The last condition holds exactly
when at least one generator has nonzero residue.
This decides exact equality even for a singular
or nonreduced $T$.

For two neighborhoods pass to their tensor
product over $S$, with the residue map
$x\otimes y\mapsto\rho(x)\rho'(y)$.
It is again a finite presented \'etale
neighborhood at that residue point, and its
pointed local henselization is $R$.
The same annihilator test decides equality
of the two names and permits arithmetic in
that common presentation.  The residue of
each computed element is explicit.
A unit has nonzero residue; adjoining its
inverse gives a pointed \'etale localization
and hence an exact inverse name.
Every $c\in\kappa$ has a lift in $S$.
Enumeration of $S$ with the decidable test
$\rho_S(s)=c$ finds such a lift in finite time.
No coefficient-field section of $R\to\kappa$
is required.

Raw names form a decidable subset of natural
number codes.  Internal tensor/localization
names can be returned to the standard list
by enumerating raw names and testing equality;
coverage proves termination.  Retaining a
valid code precisely when no smaller valid
code names the same element gives exact
canonical finite names, using finitely many
predecessor tests.

Finally, put every entry of an input matrix
in one tensor neighborhood $T$.
Compute generators of its actual kernel over
$T$.  Flatness of $T\to R$ identifies their
base change with the full kernel over $R$.
This proves the claimed kernel algorithm.
The ring $R$ is Noetherian henselian local
with residue $\kappa$, by the Noetherian
henselization theorem
\cite[Lemma 15.46.3, Tag 06LJ]{C-StacksHenselianEffective}.
All its asserted access consists of finite
computations or the explicitly justified
positive searches above.
\end{proof}

\begin{theorem}[Proved: full ordinary classification over effective $C_1$ henselizations]
\label{C:C1-henselization-WG}
For every $R$ in Lemma~\ref{C:C1-henselization-access}
and every $m\geq0$, indecomposable finite projective
rational $O_{m,R}$-modules have a decidable literal
canonical index set $\mathscr I_R\subseteq\mathbf N$.
Isomorphism, indecomposability, self-duality,
nonsingular invariant form sign, actual isometry,
and equality in the split group are decidable
by ordinary terminating algorithms.

Each self-dual index has a computed reference form
$q_i$ of sign $\varepsilon_i$, a residue
endomorphism field $K_i$ given as a finite tower
over $\kappa$, and a computed adjoint involution
$\sigma_i$.  With symmetric, alternating, unitary
and paired types as in (C.3.Q2), actual computed
generators give
\[
 \begin{split}
 WG_{\rm sp}(BO_{m,R})\simeq{}&
 \bigoplus_{i\in\mathscr I_{\rm s}}
       \bigl(\mathbf Z\oplus K_i^\times/(K_i^\times)^2\bigr)\\
 &\quad\oplus
 \bigoplus_{i\in\mathscr I_{\rm a}\cup
                  \mathscr I_{\rm u}\cup
                  \mathscr I_{\rm h}}\mathbf Z.
 \end{split}                                                \tag{C.3.CH3}
\]
Square classes have decidable literal least-code
representatives.  All products and every specified
exterior-power coefficient of finite virtual inputs
are computed in finite time.  The result includes
imperfect residue fields and inseparable finite
residue endomorphism extensions.
\end{theorem}
\begin{proof}
Finite projectives over $R$ are free.
Given an idempotent presentation, independent
columns of its residue image give a free basis
of its image: a unit minor proves independence,
and Nakayama proves generation.
Coaction matrices in the finite free integral
coefficient modules $C_d$ therefore suffice.
The preceding lemma computes the full actual
Hom modules and the invariant symmetric and
alternating matrix modules as kernels of
finite linear coefficient systems.

The residue field $\kappa$ is infinite and
$C_1$, by Lemma~\ref{C:C1_field_forms}.
It has the required factorization procedures
over every finite tower, since concatenating
with its supplied tower over $k$ gives one
of the input towers.
If $H_1,\ldots,H_s$ generate a module of
$n$-square matrices defined by one of the
linear systems, its actual residue image
is the $\kappa$-span of their residues.
Every residue scalar has a computable lift.
The determinant of $\sum_jc_j\overline H_j$
has degree at most $n$ in each variable.
Evaluate it on a grid of $n+1$ distinct
enumerated elements of $\kappa$ in each
variable.  Repeated one-variable interpolation
shows that the grid is identically zero
exactly when the determinant polynomial is.
Thus this finite test decides existence of
an invertible actual solution, and lifting
a successful coefficient vector constructs
one.  This decides module isomorphism and
nonsingular forms of each sign.  Rank zero
is treated separately.

\emph{Residual algebras and actual decomposition.}
Let $E=\operatorname{End}_{R,O_m}(M)$, computed
by its finite $R$-generators, and let $B$ be
their image algebra in
$\operatorname{End}_{\kappa}(M/\mathfrak m_RM)$.
Linear algebra over $\kappa$ computes a basis,
identity and multiplication constants for $B$.
We spell out the terminating finite-algebra
certificate used in
Theorem~\ref{C:artinian_effective_imperfect_C1}.
If $h=\dim_\kappa B$, enumerate finite products
of matrix algebras over finite field towers
over $\kappa$ with total dimension at most $h$.
Factorization at each tower stage verifies
the field presentations.  Enumerate the
$\kappa$-linear maps from $B$ to each product.
Retain a map precisely when it is unital,
multiplicative, surjective and its kernel
$J$ satisfies $J^{h+1}=0$.
All these are finite exact tests.
A retained kernel is exactly
$\operatorname{rad}B$.

A certificate exists.  The semisimple
quotient is a product of matrix algebras
over finite-dimensional division algebras.
Each center is a finite extension of the
$C_1$ field $\kappa$, hence is $C_1$ and
has trivial Brauer group by
Lemma~\ref{C:C1_field_forms}.
Every division algebra is therefore its
center, and each center has a finite tower
presentation, without a separability or
primitive-element assumption.
The radical is nilpotent of exponent at
most $h+1$.  The true quotient therefore
occurs in the enumeration and proves its
termination.

The kernel of $E\to B$ lies in
$\operatorname{rad}E$: if $u$ has zero
residue, then $1-vu$ has unit determinant
for every $v\in E$, and its matrix inverse
is again equivariant.  Thus
\[
 E/\operatorname{rad}E
      \simeq B/\operatorname{rad}B.                       \tag{C.3.CH4}
\]
The ring $E$ is finite over the Noetherian
henselian local ring $R$, so
Lemma~\ref{C:henselian-finite-algebra-lifting}
makes it semiperfect.  The module $M$ is
indecomposable exactly when the computed
quotient in (C.3.CH4) is one field factor.
Otherwise choose a nontrivial matrix-block
idempotent in that quotient and lift it
to $B$ using a polynomial which is zero
modulo $T^{h+1}$ and one modulo
$(T-1)^{h+1}$.
Enumerate actual elements of $E$ until an
idempotent with that prescribed residue
in $B$ occurs.  The lifting lemma, applied
to the kernel of $E\to B$, proves existence;
exact equality tests each candidate.
This positive search therefore terminates.
Split its free image and complement by
residue columns and repeat.  Positive
ranks decrease, giving a complete
Krull--Schmidt decomposition.

\emph{Canonical indices and form parameters.}
Code finite coaction matrices by the finite
ring names from the lemma.  Coaction
identities, minimum coefficient degree,
and indecomposability are decidable.
Retain an indecomposable code precisely
when no smaller valid code is isomorphic
to it.  There are finitely many such
tests, and each terminates by the
determinant procedure.  Every class
has a least code; the same predecessor
procedure computes it from any input.
The antipode and this canonicalization
compute duality and the paired indices.

For a self-dual indecomposable, at least
one nonsingular sign exists.  If
$g:M\to M^*$ is an isomorphism, the
endomorphisms $g^{-1}(g+g^{\mathsf t})$
and $g^{-1}(g-g^{\mathsf t})$ sum to
$2$.  They cannot both lie in the
radical of the local endomorphism ring.
Thus one of the two forms is invertible.
The determinant tests decide both signs;
choose symmetric when possible, otherwise
alternating, and choose the least valid
form code of that sign, bounded by one
already constructed successful code.
This supplies $q_i$ and $\varepsilon_i$.
The computed quotient in (C.3.CH4) is
now the finite tower field $K_i$.
Matrix operations compute the adjoint
and its induced involution $\sigma_i$.
It fixes $\kappa$, since the adjoint
fixes all scalar endomorphisms over $R$.

The henselian hermitian lifting lemma
identifies actual form isometry with
the isometry of residue multiplicity
forms.  Transfer between distinct
indecomposable types is exactly the
semisimple-block argument established
there.  The finite extensions $K_i/\kappa$
are $C_1$, and when $\sigma_i\ne1$
its fixed field is a finite $C_1$
extension of $\kappa$ as well.
Lemma~\ref{C:C1_field_forms} therefore
gives rank and determinant square class
in symmetric type, half-rank in
alternating type, rank in unitary
type, and common multiplicity in a
dual pair.  This proves (C.3.CH3).
It does not use separability of
$K_i/\kappa$.

Square-class equality in $K_i$ is
decided by factoring $X^2-a/b$ in
its supplied tower.  Among the
nonzero field codes no larger than
an input code choose the least
square-equivalent one.  This finite
procedure gives literal canonical
square classes and their computable
products and inverses.

\emph{Actual forms, coordinates and operations.}
For a symmetric index use $[q_i]$
for its integer generator.
For a square-class scalar $a\in K_i^\times$,
lift it through the computed residue
algebra and the actual Hom generators,
using the residue lifts from the lemma.
If the resulting endomorphism is $b$,
replace it by $b'=(b+b^*)/2$.
Its residue is $a$, so it is a unit,
and $[q_i b']-[q_i]$ represents the
rank-zero square-class coordinate.
In alternating and paired types
use $H(M_i)$.
In unitary type choose nonzero
$a\in K_i$ with
$\sigma_i(a)=\varepsilon_i a$:
use $a=1$ for positive sign, or
$a=c-\sigma_i(c)\ne0$ for negative
sign.  A $\kappa$-basis finds such
a $c$.  Lift and average as
$(b+\varepsilon_i b^*)/2$ to get
an actual symmetric generator.
Every finite group coordinate
therefore has a computable finite
combination of actual representatives.

For an input actual form, compute
its underlying Krull--Schmidt
decomposition and the isomorphisms
to the canonical representatives.
If $B_{uv}:M_i\to M_i^*$ are the
blocks on a self-dual type, its
residual multiplicity matrix is
\[
       H_{uv}=\pi_i(q_i^{-1}B_{uv})\in K_i.              \tag{C.3.CH5}
\]
It is nonsingular: reducing the
global form isomorphism to the
semisimple endomorphism category
kills only radical off-diagonal
maps between different types.
The fixed self-dual diagonal
block must remain invertible.
Its determinant square class
or its multiplicity dimension
gives the coordinate just proved;
paired types give their common
multiplicity.  Thus all actual
coordinates are computed.
They decide actual form isometry,
because the residue-form monoids
have exactly these invariants
and henselian lifting is exact.
If a matrix realizing a positive
isometry decision is requested,
enumerate actual intertwiner
matrices and test the Gram and
inverse identities until one
appears.  The decision proves
existence, hence termination.
Cancellation of the block monoids
also proves equality in their
group completions is exactly
equality of (C.3.CH3).

Tensor products use coefficient
products and tensor Gram matrices;
exterior powers use minors for
both coactions and forms.
Their coefficient degrees are
bounded by $d+d'$ and $rd$.
Apply the full decomposition
and coordinate procedures above.
Extend multiplication bilinearly
and expand
$\lambda_t(-x)=\lambda_t(x)^{-1}$
to each requested finite order.
Every subroutine terminates, so
this gives all asserted operations.
\end{proof}

\begin{corollary}[Proved: positive-dimensional function-field coefficient bases]
\label{C:function-field-henselization-WG}
Fix $p>2$, and let $k$ be any supplied finite
field tower over $\overline{\mathbf F}_p(t)$.
For every finitely presented $k$-algebra $S$
and supplied surjection to a finite tower
field $\kappa/k$, all conclusions of
Theorem~\ref{C:C1-henselization-WG} hold
for $(S_{\ker\rho_S})^h$.
In particular this gives ordinary complete
classification and operations over
\[
 (k[x]_{(x)})^h,\qquad
 \bigl((k[x,y]/(xy))_{(x,y)}\bigr)^h,
 \qquad
 \bigl((k[x,y]/(x^2))_{(x,y)}\bigr)^h.                   \tag{C.3.CH6}
\]
The latter two examples respectively include
a reduced singular local base and a
nonreduced local base of positive dimension.
No unspecified factorization oracle remains.
\end{corollary}
\begin{proof}
Theorem~\ref{C:function-field-factorization}
gives exact arithmetic and factorization
over $k$ and every further finite tower.
Tsen's theorem and the finite-extension
norm argument, as applied in
Theorem~\ref{C:function-field-WG},
make $k$ a $C_1$ field, including
inseparable displayed towers.
The field is countable and infinite.
Thus every hypothesis of the preceding
theorem is supplied by a proved ordinary
algorithm for this actual family.

In (C.3.CH6) the residue map evaluates
the displayed variables at zero, with
residue field $k$.
The first local ring is a discrete
valuation ring.  The nodal ring has
dimension one and cotangent dimension
two, and is reduced; henselization
preserves reducedness and regularity
\cite[Lemmas 15.46.4 and 15.46.10,
Tag 07QL]{C-StacksHenselianEffective}.
Its dimension remains one.
In the last ring $x$ is a nonzero
square-zero element, which remains
nonzero under the faithfully flat
henselization.  Its dimension is
one, since its nilradical quotient
is $k[y]_{(y)}$.
Dimension is preserved under
henselization by Lemma 15.46.7
of the same source.
These verify the stated examples
and complete the proof.
\end{proof}
